% ============================================================
%  SY05-28 曲线的镜像及其 Fukaya 范畴 —— 中文译注本
%  preamble: 版式、定理环境、宏、TikZ 图样
% ============================================================
\documentclass[11pt,a4paper,fontset=windows]{ctexart}

\usepackage[a4paper,margin=2.6cm,top=2.9cm,bottom=2.7cm]{geometry}
\usepackage{amsmath,amssymb,amsthm,mathtools,mathrsfs}
\usepackage{graphicx}
\usepackage{array,booktabs,longtable,multirow}
\usepackage{caption}
\usepackage{enumitem}
\usepackage{xcolor}
\usepackage{fancyhdr}
\usepackage{tikz}
\usetikzlibrary{arrows.meta,calc,decorations.pathmorphing,decorations.markings,
  patterns,positioning,shapes.geometric,backgrounds}
\usepackage[colorlinks=true,linkcolor=blue!55!black,citecolor=blue!55!black,
  urlcolor=blue!55!black,bookmarksnumbered=true,
  pdftitle={曲线的镜像及其 Fukaya 范畴},pdfauthor={Denis Auroux（著）；数译（译注）}]{hyperref}

% ---------------- 定理环境（共享计数，按节编号） ----------------
\theoremstyle{plain}
\newtheorem{theo}{定理}[section]
\newtheorem{lemma}[theo]{引理}
\newtheorem{prop}[theo]{命题}
\newtheorem{cor}[theo]{推论}
\theoremstyle{definition}
\newtheorem{defi}[theo]{定义}
\newtheorem{conj}[theo]{猜想}
\newtheorem{example}[theo]{例}
\newtheorem{cond}[theo]{条件}
\theoremstyle{remark}
\newtheorem{remark}[theo]{注}
\newtheorem{property}[theo]{性质}
\numberwithin{equation}{section}

% ---------------- 版式 ----------------
\setlength{\parindent}{2em}
\linespread{1.28}
\setlist{nosep,leftmargin=2.2em}
\renewcommand{\arraystretch}{1.15}
\captionsetup{font=small,labelsep=quad}

\pagestyle{fancy}
\fancyhf{}
\fancyhead[L]{\small\kaishu 曲线的镜像及其 Fukaya 范畴}
\fancyhead[R]{\small\thepage}
\renewcommand{\headrulewidth}{0.4pt}
\fancypagestyle{plain}{\fancyhf{}\fancyhead[R]{\small\thepage}%
  \renewcommand{\headrulewidth}{0pt}}

\ctexset{
  section/format={\Large\bfseries\heiti},
  subsection/format={\large\bfseries\heiti},
  subsubsection/format={\normalsize\bfseries\heiti},
}

% ---------------- 宏 ----------------
\newcommand{\Z}{\mathbb{Z}}
\newcommand{\R}{\mathbb{R}}
\newcommand{\C}{\mathbb{C}}
\newcommand{\Q}{\mathbb{Q}}
\newcommand{\N}{\mathbb{N}}
\newcommand{\K}{\mathbb{K}}
\newcommand{\CP}{\mathbb{CP}}
\newcommand{\A}{\mathbb{A}}
\newcommand{\PP}{\mathbb{P}}
\newcommand{\OO}{\mathcal{O}}
\newcommand{\LL}{\mathcal{L}}
\newcommand{\FF}{\mathcal{F}}
\newcommand{\GG}{\mathcal{G}}
\newcommand{\EE}{\mathcal{E}}
\newcommand{\crit}{\operatorname{crit}}
\newcommand{\Trop}{\operatorname{Trop}}
\newcommand{\Newt}{\operatorname{Newt}}
\newcommand{\Conv}{\operatorname{Conv}}
\newcommand{\Tot}{\operatorname{Tot}}
\newcommand{\Coh}{\mathrm{Coh}}
\newcommand{\Perf}{\mathrm{Perf}}
\newcommand{\Ext}{\operatorname{Ext}}
\newcommand{\Hom}{\operatorname{Hom}}
\newcommand{\End}{\operatorname{End}}
\newcommand{\Aut}{\operatorname{Aut}}
\newcommand{\val}{\operatorname{val}}
\newcommand{\Spec}{\operatorname{Spec}}
\newcommand{\Log}{\operatorname{Log}}
\newcommand{\id}{\mathrm{id}}
\newcommand{\db}{D^{b}}
\newcommand{\dsg}{D^{b}_{\mathrm{sg}}}
\newcommand{\Tw}{\mathrm{Tw}}
\newcommand{\Fwr}{F^{\mathrm{wr}}}
\newcommand{\hms}{\mathrm{HMS}}
\newcommand{\loc}{\mathrm{loc}}
\newcommand{\ol}{\overline}
\newcommand{\wt}{\widetilde}
\newcommand{\sgn}{\operatorname{sgn}}
\newcommand{\cyc}{\mathrm{cyc}}
\newcommand{\ev}{\mathrm{ev}}
\newcommand{\bdl}{\mathrm{bdry}}
\newcommand{\xto}{\xrightarrow}
\newcommand{\lto}{\longrightarrow}
\newcommand{\mto}{\mapsto}

% 封面与图注用的短书名
\newcommand{\booktitle}{曲线的镜像及其 Fukaya 范畴}
\newcommand{\booksubtitle}{三场报告与一份导引}

% 原稿页码引用
\newcommand{\origpage}[1]{\textbf{（原稿第 #1 页）}}

\setcounter{tocdepth}{2}
